\chapter{Chapter06 — Restricted Weight Algorithm}
%%% generated by the blueprint pipeline from FABL.Chapter06.LearningAndTesting.RestrictedWeightAlgorithm
\section{Overview}

Declarations of the Lean module \texttt{FABL.Chapter06.LearningAndTesting.RestrictedWeightAlgorithm}.

\section{Declarations}

\begin{definition}[Restricted-weight input]
\lean{FABL.RestrictedWeightInput}
\statementsource{Analysis-of-Boolean-Functions-by-Ryan-ODonnell}{pdf-ffbe9c7b3306-p158-b003}
\label{FABL.RestrictedWeightInput}
\leanok
\uses{FABL.SmallBiasInput}
A restricted-weight input consists of a small-bias input and a natural-number upper
bound for the Fourier $1$-norm that is at least $1$.
\end{definition}

\begin{definition}[Encoded bias parameter]
\lean{FABL.RestrictedWeightInput.epsilon}
\statementsource{Analysis-of-Boolean-Functions-by-Ryan-ODonnell}{pdf-ffbe9c7b3306-p158-b003}
\label{FABL.RestrictedWeightInput.epsilon}
\leanok
\uses{FABL.RestrictedWeightInput, FABL.SmallBiasInput.epsilon}
For a restricted-weight input whose bias has numerator $a$ and denominator $b$, its real
accuracy parameter is $a/b$.
\end{definition}

\begin{definition}[Quartered bias input]
\lean{FABL.RestrictedWeightInput.quarterBias}
\label{FABL.RestrictedWeightInput.quarterBias}
\leanok
\uses{FABL.RestrictedWeightInput, FABL.SmallBiasInput}
Given a restricted-weight input, define its quarter-bias input to have the same output
dimension and bias numerator, and a bias denominator equal to four times the original
denominator. The resulting input has positive dimension, positive bias numerator, and
bias at most one half.
\end{definition}

\begin{theorem}[Accuracy of the quarter bias]
\lean{FABL.RestrictedWeightInput.quarterBias_epsilon}
\label{FABL.RestrictedWeightInput.quarterBias_epsilon}
\leanok
\uses{FABL.RestrictedWeightInput, FABL.RestrictedWeightInput.epsilon, FABL.RestrictedWeightInput.quarterBias, FABL.SmallBiasInput.epsilon}
\difficulty{4}
For a restricted-weight input, the accuracy parameter of its quarter-bias input equals
one quarter of the input's accuracy parameter.
\end{theorem}

\begin{definition}[Inner small-bias Fourier input]
\lean{FABL.RestrictedWeightInput.innerInput}
\statementsource{Analysis-of-Boolean-Functions-by-Ryan-ODonnell}{pdf-ffbe9c7b3306-p158-b002}
\label{FABL.RestrictedWeightInput.innerInput}
\leanok
\uses{FABL.RestrictedWeightInput, FABL.RestrictedWeightInput.quarterBias, FABL.SmallBiasFourierInput}
For a restricted-weight input, form a small-bias Fourier input using the input's
quarter-bias parameters and retaining its Fourier bound.
\end{definition}

\begin{definition}[Outer small-bias Fourier input]
\lean{FABL.RestrictedWeightInput.outerInput}
\statementsource{Analysis-of-Boolean-Functions-by-Ryan-ODonnell}{pdf-ffbe9c7b3306-p158-b002}
\label{FABL.RestrictedWeightInput.outerInput}
\leanok
\uses{FABL.RestrictedWeightInput, FABL.RestrictedWeightInput.quarterBias, FABL.SmallBiasFourierInput}
For a restricted-weight input, form a small-bias Fourier input using the input's
quarter-bias parameters and the square of its Fourier bound.
\end{definition}

\begin{theorem}[Inner input bias ratio]
\lean{FABL.RestrictedWeightInput.innerInput_biasParameter}
\label{FABL.RestrictedWeightInput.innerInput_biasParameter}
\leanok
\uses{FABL.RestrictedWeightInput, FABL.RestrictedWeightInput.epsilon, FABL.RestrictedWeightInput.innerInput, FABL.RestrictedWeightInput.quarterBias_epsilon, FABL.SmallBiasFourierInput.epsilon}
\difficulty{3}
For a restricted-weight input, the inner input's accuracy parameter divided by its
Fourier bound equals the input's accuracy parameter divided by four times its Fourier
bound.
\end{theorem}

\begin{theorem}[Outer input bias ratio]
\lean{FABL.RestrictedWeightInput.outerInput_biasParameter}
\label{FABL.RestrictedWeightInput.outerInput_biasParameter}
\leanok
\uses{FABL.RestrictedWeightInput, FABL.RestrictedWeightInput.epsilon, FABL.RestrictedWeightInput.outerInput, FABL.RestrictedWeightInput.quarterBias_epsilon, FABL.SmallBiasFourierInput.epsilon}
\difficulty{3}
For a restricted-weight input, the outer input's accuracy parameter divided by its
Fourier bound equals the input's accuracy parameter divided by four times the square of
its Fourier bound.
\end{theorem}

\begin{definition}[Inner sample count]
\lean{FABL.RestrictedWeightInput.innerCount}
\label{FABL.RestrictedWeightInput.innerCount}
\leanok
\uses{FABL.RestrictedWeightInput, FABL.RestrictedWeightInput.innerInput, FABL.SmallBiasFourierInput.sampleCount}
For a restricted weight input, the inner sample count is the sample count of its
associated inner input.
\end{definition}

\begin{definition}[Outer sample count]
\lean{FABL.RestrictedWeightInput.outerCount}
\label{FABL.RestrictedWeightInput.outerCount}
\leanok
\uses{FABL.RestrictedWeightInput, FABL.RestrictedWeightInput.outerInput, FABL.SmallBiasFourierInput.sampleCount}
For a restricted weight input, the outer sample count is the sample count of its
associated outer input.
\end{definition}

\begin{theorem}[Positivity of the Inner Sample Count]
\lean{FABL.RestrictedWeightInput.innerCount_pos}
\label{FABL.RestrictedWeightInput.innerCount_pos}
\leanok
\uses{FABL.RestrictedWeightInput, FABL.RestrictedWeightInput.innerCount, FABL.RestrictedWeightInput.innerInput, FABL.SmallBiasFourierInput.sampleCount_pos}
\difficulty{1}
For every restricted-weight input—consisting of a positive-dimensional small-bias input
with positive numerator and bias at most one half, together with an integral Fourier
bound of at least one—the associated inner sample count is positive.
\end{theorem}

\begin{theorem}[Positivity of the outer sample count]
\lean{FABL.RestrictedWeightInput.outerCount_pos}
\label{FABL.RestrictedWeightInput.outerCount_pos}
\leanok
\uses{FABL.RestrictedWeightInput, FABL.RestrictedWeightInput.outerCount, FABL.RestrictedWeightInput.outerInput, FABL.SmallBiasFourierInput.sampleCount_pos}
\difficulty{1}
For every restricted-weight input, the sample count associated with its outer input is
positive.
\end{theorem}

\begin{definition}[Free sign cube equivalence]
\lean{FABL.freeSignCubeEquiv}
\label{FABL.freeSignCubeEquiv}
\leanok
\uses{FABL.FreeSignCube}
Let $J$ be a finite subset of $\\{0,\ldots,n-1\}$. The set of sign-valued functions on
$J$ is equivalent to the sign cube of dimension $|J|$, by reindexing coordinates using
the increasing enumeration of $J$.
\end{definition}

\begin{definition}[Free frequency reindexing]
\lean{FABL.freeFrequencyReindex}
\label{FABL.freeFrequencyReindex}
\leanok
Let $J$ be a finite subset of $Fin(n)$, and let $S$ be a finite subset of $J$.
Reindexing the elements of $S$ through the increasing enumeration of $J$ gives a finite
subset of $Fin(|J|)$.
\end{definition}

\begin{theorem}[Monomial under frequency reindexing]
\lean{FABL.monomial_freeFrequencyReindex}
\label{FABL.monomial_freeFrequencyReindex}
\leanok
\uses{FABL.freeFrequencyReindex, FABL.freeSignCubeEquiv, FABL.indexedMonomial, FABL.monomial}
\difficulty{3}
Let $J$ be a finite set of coordinates among $n$, let $S$ be a subset of $J$, and let
$y$ be a sign assignment to the $|J|$ positions. Reindex the elements of $S$ by their
positions in the enumeration of $J$, and transfer $y$ to an assignment on $J$ using that
enumeration. The monomial on the reindexed set evaluated at $y$ equals the monomial on
$S$ evaluated at the transferred assignment.
\end{theorem}

\begin{definition}[Fixed-coordinate projection]
\lean{FABL.fixedCoordinateProjectionLinear}
\label{FABL.fixedCoordinateProjectionLinear}
\leanok
\uses{FABL.FixedIndex, FABL.F₂Cube, FABL.𝔽₂}
Let $J$ be a finite subset of the coordinates $Fin(n)$. The fixed-coordinate projection
is the linear map from $Fin(n)\to\mathbb{F}_2$ to $Fin(|Fin(n)\setminus
J|)\to\mathbb{F}_2$ that records, in reindexed order, the coordinates outside $J$.
\end{definition}

\begin{theorem}[Coordinate projection formula]
\lean{FABL.fixedCoordinateProjectionLinear_apply}
\label{FABL.fixedCoordinateProjectionLinear_apply}
\leanok
\uses{FABL.FixedIndex, FABL.F₂Cube, FABL.fixedCoordinateProjectionLinear}
\difficulty{1}
Let $n$ be a natural number, let $J$ be a set of indices in $in n$, and let $x$ be a
function from $in n$ to the field with two elements. For any position $q$ among the
indices outside $J$, the coordinate of the projection of $x$ at $q$ equals the
coordinate of $x$ indexed by the corresponding element of $in n \setminus J$.
\end{theorem}

\begin{definition}[Fixed-coordinate extension]
\lean{FABL.fixedCoordinateExtension}
\label{FABL.fixedCoordinateExtension}
\leanok
\uses{FABL.FixedIndex, FABL.F₂Cube}
Let $J$ be a finite subset of the coordinates $Fin(n)$. Extend a function from the
coordinates outside $J$ to all coordinates by retaining its values outside $J$ and
assigning $0$ on $J$, with the outside coordinates reindexed in increasing order.
\end{definition}

\begin{theorem}[Projection of a coordinate extension]
\lean{FABL.fixedCoordinateProjectionLinear_fixedCoordinateExtension}
\label{FABL.fixedCoordinateProjectionLinear_fixedCoordinateExtension}
\leanok
\uses{FABL.FixedIndex, FABL.F₂Cube, FABL.fixedCoordinateExtension, FABL.fixedCoordinateProjectionLinear}
\difficulty{3}
Let $n$ be a natural number, let $J$ be a set of indices in $in n$, and let $z$ assign
a value in the field with two elements to each index outside $J$. Extend $z$ to all
indices in $in n$ by assigning zero on $J$. Projecting this extension onto the indices
outside $J$ recovers $z$.
\end{theorem}

\begin{definition}[Fixed-coordinate frequency lift]
\lean{FABL.fixedCoordinateFrequencyLift}
\label{FABL.fixedCoordinateFrequencyLift}
\leanok
\uses{FABL.FixedIndex, FABL.F₂Cube, FABL.fixedCoordinateProjectionLinear, FABL.𝔽₂}
Let $J$ be a finite subset of the coordinates $Fin(n)$. The fixed-coordinate frequency
lift sends a function $gamma$ on the reindexed coordinates outside $J$ to the unique
coefficient vector on all $n$ coordinates whose dot product with any vector over
$mathbb{F}_2$ equals the dot product of $gamma$ with that vector's restriction to the
coordinates outside $J$.
\end{definition}

\begin{theorem}[Dot product of a fixed-coordinate frequency lift]
\lean{FABL.f₂DotProduct_fixedCoordinateFrequencyLift}
\leanok
\uses{FABL.FixedIndex, FABL.F₂Cube, FABL.fixedCoordinateFrequencyLift, FABL.fixedCoordinateProjectionLinear, FABL.f₂DotProduct, FABL.𝔽₂}
\difficulty{4}
Let $n$ be a nonnegative integer, let $J$ be a subset of its coordinate indices, and let
$\gamma$ be a binary vector indexed by the complement of $J$, in the enumeration used to
order that complement. For every binary vector $x$ of length $n$, the dot product of the
fixed-coordinate frequency lift of $\gamma$ with $x$ equals the dot product of $\gamma$
with the restriction of $x$ to the complement of $J$.
\end{theorem}

\begin{theorem}[Walsh characters under projection away from fixed coordinates]
\lean{FABL.vectorWalshCharacter_fixedCoordinateProjection}
\label{FABL.vectorWalshCharacter_fixedCoordinateProjection}
\leanok
\uses{FABL.FixedIndex, FABL.F₂Cube, FABL.fixedCoordinateFrequencyLift, FABL.fixedCoordinateProjectionLinear, FABL.f₂DotProduct_fixedCoordinateFrequencyLift, FABL.vectorWalshCharacter, FABL.vectorWalshCharacter_apply}
\difficulty{1}
Let $n$ be a nonnegative integer, let $J\subseteq\{0,\ldots,n-1\}$, let $\gamma$ be a
binary vector indexed by the complement of $J$, and let $x\in\mathbb F_2^n$. Project $x$
onto the coordinates outside $J$, in the ordering used to identify that complement with
a finite coordinate set, and lift $\gamma$ to a vector in $\mathbb F_2^n$ by extending
its corresponding frequency vector by zero on $J$. The Walsh character indexed by
$\gamma$ evaluated at the projection of $x$ equals the Walsh character indexed by this
lifted vector evaluated at $x$.
\end{theorem}

\begin{theorem}[Nonvanishing of the lifted frequency]
\lean{FABL.fixedCoordinateFrequencyLift_ne_zero}
\label{FABL.fixedCoordinateFrequencyLift_ne_zero}
\leanok
\uses{FABL.FixedIndex, FABL.F₂Cube, FABL.fixedCoordinateExtension, FABL.fixedCoordinateFrequencyLift, FABL.fixedCoordinateProjectionLinear_fixedCoordinateExtension, FABL.f₂DotProduct, FABL.f₂DotProduct_fixedCoordinateFrequencyLift, FABL.𝔽₂}
\difficulty{4}
Let $n\in\mathbb{N}$, let $J$ be a subset of the $n$ coordinate indices, and let
$\gamma$ be a nonzero $\mathbb{F}_2$-valued frequency on the indices outside $J$. Extend
$\gamma$ to all $n$ coordinates by assigning frequency zero on $J$. The resulting
frequency is nonzero.
\end{theorem}

\begin{definition}[Restricted-weight free samples]
\lean{FABL.restrictedWeightFreeSample}
\statementsource{Analysis-of-Boolean-Functions-by-Ryan-ODonnell}{pdf-ffbe9c7b3306-p158-b002}
\label{FABL.restrictedWeightFreeSample}
\leanok
\uses{FABL.F₂Cube, FABL.RestrictedWeightInput, FABL.RestrictedWeightInput.innerInput, FABL.SmallBiasFourierInput.sample, FABL.SmallBiasFourierInput.sampleCount, FABL.coordinateProjectionLinear}
Let $n$ be the dimension of a restricted-weight input and let
$J\subseteq\mathrm{Fin}(n)$. For each index in the sample set determined by the inner
input, take the corresponding generator sample and project it onto the coordinates in
$J$, obtaining an element of $\mathbb{F}_2^J$.
\end{definition}

\begin{definition}[Restricted-weight fixed-coordinate samples]
\lean{FABL.restrictedWeightFixedSample}
\statementsource{Analysis-of-Boolean-Functions-by-Ryan-ODonnell}{pdf-ffbe9c7b3306-p158-b002}
\label{FABL.restrictedWeightFixedSample}
\leanok
\uses{FABL.FixedIndex, FABL.F₂Cube, FABL.RestrictedWeightInput, FABL.RestrictedWeightInput.outerInput, FABL.SmallBiasFourierInput.sample, FABL.SmallBiasFourierInput.sampleCount, FABL.fixedCoordinateProjectionLinear}
Let $n$ be the dimension of a restricted-weight input and let
$J\subseteq\mathrm{Fin}(n)$. For each index in the sample set determined by the outer
input, take the corresponding generator sample and project it onto the indices outside
$J$, obtaining an element of $\mathbb{F}_2^{\mathrm{Fin}(n)\setminus J}$.
\end{definition}

\begin{definition}[Restricted-weight free sign assignments]
\lean{FABL.restrictedWeightFreeAssignment}
\statementsource{Analysis-of-Boolean-Functions-by-Ryan-ODonnell}{pdf-ffbe9c7b3306-p158-b002}
\label{FABL.restrictedWeightFreeAssignment}
\leanok
\uses{FABL.FreeSignCube, FABL.RestrictedWeightInput, FABL.RestrictedWeightInput.innerInput, FABL.SmallBiasFourierInput.sampleCount, FABL.binaryCubeSignEquiv, FABL.freeSignCubeEquiv, FABL.restrictedWeightFreeSample}
Let $n$ be the dimension of a restricted-weight input and let
$J\subseteq\mathrm{Fin}(n)$. For each index in the sample set determined by the inner
input, convert the corresponding projected binary sample into a sign assignment on $J$.
\end{definition}

\begin{definition}[Restricted weight fixed assignment]
\lean{FABL.restrictedWeightFixedAssignment}
\statementsource{Analysis-of-Boolean-Functions-by-Ryan-ODonnell}{pdf-ffbe9c7b3306-p158-b002}
\label{FABL.restrictedWeightFixedAssignment}
\leanok
\uses{FABL.FixedIndex, FABL.FixedSignCube, FABL.RestrictedWeightInput, FABL.RestrictedWeightInput.outerInput, FABL.SmallBiasFourierInput.sampleCount, FABL.binaryCubeSignEquiv, FABL.fixedSignCubeEquiv, FABL.restrictedWeightFixedSample}
Let $n>0$, and let $a,b$ be positive natural numbers with $2a\le b$. Given a
restricted-weight input with these parameters, a subset $J\subseteq\mathrm{Fin}(n)$, and
a sample index in the input's outer sample range, the fixed assignment is the sign
assignment on $\mathrm{Fin}(n)\setminus J$ obtained by encoding the restricted-weight
fixed sample at that index as a sign cube.
\end{definition}

\begin{theorem}[Bias of the projected sample]
\lean{FABL.restrictedWeightFreeSample_isBiased}
\label{FABL.restrictedWeightFreeSample_isBiased}
\leanok
\uses{FABL.ProbabilityDensity.IsBiased, FABL.ProbabilityDensity.expectation_uniformPushforward, FABL.ProbabilityDensity.isBiased_iff_expectation, FABL.ProbabilityDensity.uniformPushforward, FABL.RestrictedWeightInput, FABL.RestrictedWeightInput.epsilon, FABL.RestrictedWeightInput.innerInput, FABL.RestrictedWeightInput.innerInput_biasParameter, FABL.SmallBiasFourierInput.sample, FABL.SmallBiasFourierInput.sampleCount, FABL.SmallBiasFourierInput.sampleCount_pos, FABL.SmallBiasFourierInput.sampleDensity, FABL.SmallBiasFourierInput.sample_isBiased, FABL.coordinateFrequencyLift, FABL.coordinateFrequencyLift_ne_zero, FABL.restrictedWeightFreeSample, FABL.vectorWalshCharacter_coordinateProjection}
\difficulty{4}
Let $n,a,b,s$ be integers with $n>0$, $a>0$, $2a\le b$, and $s\ge1$. Form the
restricted-weight input with dimension $n$, target bias $\varepsilon=a/b$, and Fourier
bound $s$. For every set $J$ of coordinates among the $n$ coordinates, choose an index
uniformly from the construction's finite, nonempty sample set and project the resulting
sample onto $J$. The resulting distribution has every nonzero Fourier coefficient
bounded in absolute value by $\varepsilon/(4s)$.
\end{theorem}

\begin{theorem}[Bias of the projected sample]
\lean{FABL.restrictedWeightFixedSample_isBiased}
\label{FABL.restrictedWeightFixedSample_isBiased}
\leanok
\uses{FABL.ProbabilityDensity.IsBiased, FABL.ProbabilityDensity.expectation_uniformPushforward, FABL.ProbabilityDensity.isBiased_iff_expectation, FABL.ProbabilityDensity.uniformPushforward, FABL.RestrictedWeightInput, FABL.RestrictedWeightInput.epsilon, FABL.RestrictedWeightInput.outerInput, FABL.RestrictedWeightInput.outerInput_biasParameter, FABL.SmallBiasFourierInput.sample, FABL.SmallBiasFourierInput.sampleCount, FABL.SmallBiasFourierInput.sampleCount_pos, FABL.SmallBiasFourierInput.sampleDensity, FABL.SmallBiasFourierInput.sample_isBiased, FABL.fixedCoordinateFrequencyLift, FABL.fixedCoordinateFrequencyLift_ne_zero, FABL.restrictedWeightFixedSample, FABL.vectorWalshCharacter_fixedCoordinateProjection}
\difficulty{4}
Let $n,a,b,s$ be integers with $n>0$, $a>0$, $2a\le b$, and $s\ge1$. Set
$\varepsilon=a/b$, and let $J$ be any subset of the $n$ coordinates. Choose uniformly an
index from the sample specified by the restricted-weight input with dimension $n$,
target bias $\varepsilon$, and Fourier bound $s$, and project the resulting binary
vector onto the coordinates outside $J$. Every nonzero Fourier coefficient of the
resulting distribution has absolute value at most $\varepsilon/(4s^2)$.
\end{theorem}

\begin{theorem}[Fourier coeff reindexed sign restriction]
\lean{FABL.fourierCoeff_reindexedSignRestriction}
\statementsource{Analysis-of-Boolean-Functions-by-Ryan-ODonnell}{
  pdf-ffbe9c7b3306-p076-b003,
  pdf-ffbe9c7b3306-p076-b004
}
\label{FABL.fourierCoeff_reindexedSignRestriction}
\leanok
\uses{FABL.BooleanFunction, FABL.BooleanFunction.toReal, FABL.FixedSignCube, FABL.combineSignCube, FABL.fourierCoeff, FABL.freeSignCubeEquiv, FABL.indexedFourierCoeff, FABL.indexedMonomial, FABL.monomial, FABL.signRestriction}
\difficulty{4}
% dataset-pipeline: stub -- prose pending build_informal_statements.py
Statement pending regeneration.
\end{theorem}

\begin{theorem}[Inverse frequency reindexing]
\lean{FABL.map_freeFrequencyReindex_symm}
\label{FABL.map_freeFrequencyReindex_symm}
\leanok
\uses{FABL.freeFrequencyReindex}
\difficulty{1}
Let $n$ be a natural number, let $J$ be a finite set of elements of $\{0,\ldots,n-1\}$,
and let $S$ be a subset of $J$. Reindex $S$ by the increasing bijection from $J$ to
$\{0,\ldots,|J|-1\}$, then map the resulting subset back along the inverse bijection;
the result is $S$.
\end{theorem}

\begin{theorem}[Fourier coeff reindexed sign restriction free frequency reindex]
\lean{FABL.fourierCoeff_reindexedSignRestriction_freeFrequencyReindex}
\statementsource{Analysis-of-Boolean-Functions-by-Ryan-ODonnell}{
  pdf-ffbe9c7b3306-p076-b003,
  pdf-ffbe9c7b3306-p076-b004
}
\label{FABL.fourierCoeff_reindexedSignRestriction_freeFrequencyReindex}
\leanok
\uses{FABL.BooleanFunction, FABL.BooleanFunction.toReal, FABL.FixedSignCube, FABL.combineSignCube, FABL.fourierCoeff, FABL.fourierCoeff_reindexedSignRestriction, FABL.freeFrequencyReindex, FABL.freeSignCubeEquiv, FABL.map_freeFrequencyReindex_symm, FABL.restrictionFourierCoeff}
\difficulty{3}
% dataset-pipeline: stub -- prose pending build_informal_statements.py
Statement pending regeneration.
\end{theorem}

\begin{theorem}[Fourier norm of a reindexed restriction]
\lean{FABL.fourierOneNorm_reindexedSignRestriction_eq}
\label{FABL.fourierOneNorm_reindexedSignRestriction_eq}
\leanok
\uses{FABL.BooleanFunction, FABL.BooleanFunction.toReal, FABL.FixedSignCube, FABL.combineSignCube, FABL.fourierCoeff_reindexedSignRestriction, FABL.fourierOneNorm, FABL.freeSignCubeEquiv, FABL.indexedFourierCoeff, FABL.signRestriction}
\difficulty{4}
Let $n$ be a nonnegative integer, let $h$ be a sign-valued function on $\\{-1,1\\}^n$,
let $J$ be a set of coordinates, and fix a sign assignment $z$ on the coordinates
outside $J$. The Fourier one-norm of the restriction of $h$, viewed as a real-valued
function on a cube indexed by $J$, equals the sum of the absolute values of its Fourier
coefficients indexed by subsets of $J$.
\end{theorem}

\begin{theorem}[Fourier norm under restriction]
\lean{FABL.fourierOneNorm_reindexedSignRestriction_le}
\label{FABL.fourierOneNorm_reindexedSignRestriction_le}
\leanok
\uses{FABL.BooleanFunction, FABL.BooleanFunction.toReal, FABL.FixedSignCube, FABL.combineSignCube, FABL.fourierOneNorm, FABL.fourierOneNorm_reindexedSignRestriction_eq, FABL.freeSignCubeEquiv, FABL.sum_abs_indexedFourierCoeff_signRestriction_le_fourierOneNorm}
\difficulty{1}
Let $n$ be a nonnegative integer, let $h$ be a sign-valued function on $\\{-1,1\\}^n$,
let $J$ be a set of coordinates, and fix a sign assignment $z$ outside $J$. The Fourier
one-norm of the restriction of $h$ to the coordinates in $J$ is at most the Fourier
one-norm of $h$.
\end{theorem}

\begin{definition}[Restricted Fourier coefficient estimate]
\lean{FABL.restrictedFourierCoefficientEstimate}
\statementsource{Analysis-of-Boolean-Functions-by-Ryan-ODonnell}{
  pdf-ffbe9c7b3306-p080-b002,
  pdf-ffbe9c7b3306-p157-b008
}
\label{FABL.restrictedFourierCoefficientEstimate}
\leanok
\uses{FABL.BooleanFunction, FABL.BooleanFunction.toReal, FABL.FixedSignCube, FABL.RestrictedWeightInput, FABL.RestrictedWeightInput.innerInput, FABL.SmallBiasFourierInput.sampleCount, FABL.combineSignCube, FABL.indexedMonomial, FABL.restrictedWeightFreeAssignment}
Let $n>0$ and let $a,b$ be positive natural numbers with $2a\le b$. Given a Boolean
function $f:\{−1,1\}^n\to\{−1,1\}$, a subset $J\subseteq\mathrm{Fin}(n)$, a subset
$S\subseteq J$, and a sign assignment $z$ on $\mathrm{Fin}(n)\setminus J$, define the
restricted Fourier coefficient estimate to be the average over the inner sample range of
$f$ evaluated at the combined free and fixed assignments, multiplied by the character
indexed by $S$ of the free assignment.
\end{definition}

\begin{theorem}[Restricted fourier coefficient estimate eq small bias fourier estimate]
\lean{FABL.restrictedFourierCoefficientEstimate_eq_smallBiasFourierEstimate}
\statementsource{Analysis-of-Boolean-Functions-by-Ryan-ODonnell}{pdf-ffbe9c7b3306-p157-b009}
\label{FABL.restrictedFourierCoefficientEstimate_eq_smallBiasFourierEstimate}
\leanok
\uses{FABL.BooleanFunction, FABL.BooleanFunction.toReal, FABL.FixedSignCube, FABL.RestrictedWeightInput, FABL.RestrictedWeightInput.innerInput, FABL.SmallBiasFourierInput.sampleCount, FABL.combineSignCube, FABL.freeFrequencyReindex, FABL.freeSignCubeEquiv, FABL.monomial_freeFrequencyReindex, FABL.restrictedFourierCoefficientEstimate, FABL.restrictedWeightFreeSample, FABL.smallBiasFourierEstimate}
\difficulty{3}
% dataset-pipeline: stub -- prose pending build_informal_statements.py
Statement pending regeneration.
\end{theorem}

\begin{theorem}[Restricted Fourier coefficient estimation error]
\lean{FABL.abs_restrictedFourierCoefficientEstimate_sub_le}
\label{FABL.abs_restrictedFourierCoefficientEstimate_sub_le}
\leanok
\uses{FABL.BooleanFunction, FABL.BooleanFunction.toReal, FABL.FixedSignCube, FABL.RestrictedWeightInput, FABL.RestrictedWeightInput.epsilon, FABL.RestrictedWeightInput.innerInput, FABL.SmallBiasFourierInput.sampleCount, FABL.SmallBiasFourierInput.sampleCount_pos, FABL.SmallBiasInput.epsilon_pos, FABL.abs_smallBiasFourierEstimate_sub_fourierCoeff_le, FABL.combineSignCube, FABL.fourierCoeff, FABL.fourierCoeff_reindexedSignRestriction_freeFrequencyReindex, FABL.fourierOneNorm, FABL.fourierOneNorm_reindexedSignRestriction_le, FABL.freeFrequencyReindex, FABL.freeSignCubeEquiv, FABL.restrictedFourierCoefficientEstimate, FABL.restrictedFourierCoefficientEstimate_eq_smallBiasFourierEstimate, FABL.restrictedWeightFreeSample, FABL.restrictedWeightFreeSample_isBiased, FABL.restrictionFourierCoeff, FABL.smallBiasFourierEstimate}
\difficulty{4}
Let $n>0$, and let $\varepsilon=a/b$ be specified by positive integers $a,b$ with $2a\le
b$. Let $s\ge1$ be an integer, and let $f:\{−1,1\}^n\to\{−1,1\}$ have Fourier $1$-norm
at most $s$. For any coordinate set $J$, any subset $S$ of its indices, and any fixed
sign assignment $z$ on the complementary coordinates, the absolute difference between
the restricted Fourier coefficient estimate for $f$ at $(J,S,z)$ and the corresponding
restricted Fourier coefficient is at most $\varepsilon/4$.
\end{theorem}

\begin{theorem}[Absolute value of a sign]
\lean{FABL.abs_signValue_eq_one}
\label{FABL.abs_signValue_eq_one}
\leanok
\uses{FABL.Sign, FABL.signValue, FABL.signValue_neg_one, FABL.signValue_one}
\difficulty{4}
For every sign $s\in\{ -1,1\}$, its real value has absolute value $1$.
\end{theorem}

\begin{theorem}[Absolute value of an indexed sign monomial]
\lean{FABL.abs_indexedMonomial_eq_one'}
\label{FABL.abs_indexedMonomial_eq_one'}
\leanok
\uses{FABL.IndexedSignCube, FABL.indexedMonomial, FABL.indexedMonomial_sq}
\difficulty{4}
Let $\\iota$ be an index set, let $S\\subseteq\\iota$ be finite, and assign a sign to
each index in $\\iota$. The absolute value of the product of the signs indexed by $S$ is
$1$.
\end{theorem}

\begin{theorem}[Bound on a restricted Fourier coefficient]
\lean{FABL.abs_restrictionFourierCoeff_toReal_le_one}
\label{FABL.abs_restrictionFourierCoeff_toReal_le_one}
\leanok
\uses{FABL.BooleanFunction, FABL.BooleanFunction.toReal, FABL.FixedSignCube, FABL.FreeSignCube, FABL.abs_indexedMonomial_eq_one', FABL.abs_signValue_eq_one, FABL.indexedFourierCoeff, FABL.indexedMonomial, FABL.restrictionFourierCoeff, FABL.signRestriction, FABL.signRestriction_apply}
\difficulty{3}
Let $n$ be a natural number, let $f:\{−1,1\}^n\to\{−1,1\}$, and choose a coordinate set
$J$, a subset $S$ of $J$, and a fixed sign assignment on the complementary coordinates.
The absolute value of the resulting restricted Fourier coefficient is at most $1$.
\end{theorem}

\begin{definition}[Outer mean of restricted Fourier weight]
\lean{FABL.restrictedFourierWeightOuterMean}
\statementsource{Analysis-of-Boolean-Functions-by-Ryan-ODonnell}{
  pdf-ffbe9c7b3306-p084-b002,
  pdf-ffbe9c7b3306-p084-b003,
  pdf-ffbe9c7b3306-p084-b005,
  pdf-ffbe9c7b3306-p158-b002
}
\label{FABL.restrictedFourierWeightOuterMean}
\leanok
\uses{FABL.BooleanFunction, FABL.BooleanFunction.toReal, FABL.RestrictedWeightInput, FABL.RestrictedWeightInput.outerInput, FABL.SmallBiasFourierInput.sampleCount, FABL.restrictedWeightFixedAssignment, FABL.restrictionFourierCoeff}
Let $n>0$ and let $a,b$ be positive natural numbers with $2a\le b$. Given a Boolean
function $f$ on the $n$-dimensional sign cube, a subset $J\subseteq\mathrm{Fin}(n)$, and
$S\subseteq J$, define the outer mean of restricted Fourier weight as the average over
the outer sample range of the square of the restricted Fourier coefficient of $f$ on
$S$, evaluated at the fixed assignment generated by each sample.
\end{definition}

\begin{definition}[Restricted Fourier weight estimate]
\lean{FABL.restrictedFourierWeightEstimate}
\statementsource{Analysis-of-Boolean-Functions-by-Ryan-ODonnell}{
  pdf-ffbe9c7b3306-p084-b004,
  pdf-ffbe9c7b3306-p084-b005,
  pdf-ffbe9c7b3306-p158-b003
}
\label{FABL.restrictedFourierWeightEstimate}
\leanok
\uses{FABL.BooleanFunction, FABL.RestrictedWeightInput, FABL.RestrictedWeightInput.outerInput, FABL.SmallBiasFourierInput.sampleCount, FABL.restrictedFourierCoefficientEstimate, FABL.restrictedWeightFixedAssignment}
Let $n>0$ and let $a,b$ be positive natural numbers with $2a\le b$. Given a Boolean
function $f$ on the $n$-dimensional sign cube, a subset $J\subseteq\mathrm{Fin}(n)$, and
$S\subseteq J$, define the restricted Fourier weight estimate as the average over the
outer sample range of the square of the restricted Fourier coefficient estimate
evaluated at each sample's fixed assignment.
\end{definition}

\begin{theorem}[Restricted fourier weight eq mean fixed reindex]
\lean{FABL.restrictedFourierWeight_eq_mean_fixedReindex}
\statementsource{Analysis-of-Boolean-Functions-by-Ryan-ODonnell}{
  pdf-ffbe9c7b3306-p084-b002,
  pdf-ffbe9c7b3306-p084-b003,
  pdf-ffbe9c7b3306-p084-b005
}
\label{FABL.restrictedFourierWeight_eq_mean_fixedReindex}
\leanok
\uses{FABL.FixedIndex, FABL.fixedSignCubeEquiv, FABL.mean, FABL.restrictedFourierWeight, FABL.restrictedFourierWeight_equation6_6, FABL.restrictionFourierCoeff}
\difficulty{3}
% dataset-pipeline: stub -- prose pending build_informal_statements.py
Statement pending regeneration.
\end{theorem}

\begin{theorem}[Outer mean approximation of restricted Fourier weight]
\lean{FABL.abs_restrictedFourierWeightOuterMean_sub_le}
\label{FABL.abs_restrictedFourierWeightOuterMean_sub_le}
\leanok
\uses{FABL.BooleanFunction, FABL.BooleanFunction.toReal, FABL.FixedIndex, FABL.ProbabilityDensity, FABL.ProbabilityDensity.IsBiased, FABL.ProbabilityDensity.expectation_uniformPushforward, FABL.ProbabilityDensity.uniformPushforward, FABL.RestrictedWeightInput, FABL.RestrictedWeightInput.epsilon, FABL.RestrictedWeightInput.outerInput, FABL.SmallBiasFourierInput.sampleCount, FABL.SmallBiasFourierInput.sampleCount_pos, FABL.SmallBiasInput.epsilon_pos, FABL.fourierOneNorm, FABL.restrictedFourierWeight, FABL.restrictedFourierWeightOuterMean, FABL.restrictedFourierWeight_eq_mean_fixedReindex, FABL.restrictedWeightFixedAssignment, FABL.restrictedWeightFixedSample, FABL.restrictedWeightFixedSample_isBiased, FABL.restrictionFourierWeight_equation6_7}
\difficulty{4}
Let $n>0$, and let $\varepsilon=a/b$ be specified by positive integers $a,b$ with $2a\le
b$. Let $s\ge1$ be an integer, and let $f:\{−1,1\}^n\to\{−1,1\}$ have Fourier $1$-norm
at most $s$. For any coordinate set $J$ and any subset $S$ of its indices, the absolute
difference between the outer mean of the squared restricted Fourier coefficient and the
restricted Fourier weight is at most $\varepsilon/4$.
\end{theorem}

\begin{theorem}[Difference of squares estimate]
\lean{FABL.abs_sq_sub_sq_le_three_quarters}
\label{FABL.abs_sq_sub_sq_le_three_quarters}
\leanok
\difficulty{4}
Let $a,b,\varepsilon\in\mathbb R$ satisfy $0\le\varepsilon\le\tfrac12$, $|a|\le1$, and
$|b-a|\le\varepsilon/4$. Then $|b^2-a^2|\le3\varepsilon/4$.
\end{theorem}

\begin{theorem}[Restricted Fourier weight estimate error]
\lean{FABL.abs_restrictedFourierWeightEstimate_sub_outerMean_le}
\label{FABL.abs_restrictedFourierWeightEstimate_sub_outerMean_le}
\leanok
\uses{FABL.BooleanFunction, FABL.BooleanFunction.toReal, FABL.RestrictedWeightInput, FABL.RestrictedWeightInput.epsilon, FABL.RestrictedWeightInput.outerInput, FABL.SmallBiasFourierInput.sampleCount, FABL.SmallBiasFourierInput.sampleCount_pos, FABL.SmallBiasInput.epsilon_le_half, FABL.SmallBiasInput.epsilon_pos, FABL.abs_restrictedFourierCoefficientEstimate_sub_le, FABL.abs_restrictionFourierCoeff_toReal_le_one, FABL.abs_sq_sub_sq_le_three_quarters, FABL.fourierOneNorm, FABL.restrictedFourierCoefficientEstimate, FABL.restrictedFourierWeightEstimate, FABL.restrictedFourierWeightOuterMean, FABL.restrictedWeightFixedAssignment, FABL.restrictionFourierCoeff}
\difficulty{4}
Let $n>0$, and let $\varepsilon=a/b$ be specified by positive integers $a,b$ with $2a\le
b$. Let $s\ge1$ be an integer, and let $f:\{−1,1\}^n\to\{−1,1\}$ have Fourier $1$-norm
at most $s$. For any coordinate set $J$ and any subset $S$ of its indices, the absolute
difference between the restricted Fourier weight estimate and its outer mean is at most
$3\varepsilon/4$.
\end{theorem}

\begin{theorem}[Abs restricted fourier weight estimate sub le]
\lean{FABL.abs_restrictedFourierWeightEstimate_sub_le}
\statementsource{Analysis-of-Boolean-Functions-by-Ryan-ODonnell}{
  pdf-ffbe9c7b3306-p084-b004,
  pdf-ffbe9c7b3306-p084-b005,
  pdf-ffbe9c7b3306-p158-b002,
  pdf-ffbe9c7b3306-p158-b003
}
\label{FABL.abs_restrictedFourierWeightEstimate_sub_le}
\leanok
\uses{FABL.BooleanFunction, FABL.BooleanFunction.toReal, FABL.RestrictedWeightInput, FABL.RestrictedWeightInput.epsilon, FABL.abs_restrictedFourierWeightEstimate_sub_outerMean_le, FABL.abs_restrictedFourierWeightOuterMean_sub_le, FABL.fourierOneNorm, FABL.restrictedFourierWeight, FABL.restrictedFourierWeightEstimate, FABL.restrictedFourierWeightOuterMean}
\difficulty{4}
% dataset-pipeline: stub -- prose pending build_informal_statements.py
Statement pending regeneration.
\end{theorem}

\begin{definition}[Deterministic restricted weight query count]
\lean{FABL.deterministicRestrictedWeightQueryCount}
\statementsource{Analysis-of-Boolean-Functions-by-Ryan-ODonnell}{pdf-ffbe9c7b3306-p158-b003}
\label{FABL.deterministicRestrictedWeightQueryCount}
\leanok
\uses{FABL.RestrictedWeightInput, FABL.RestrictedWeightInput.innerCount, FABL.RestrictedWeightInput.outerCount}
Let $n>0$ and let $a,b$ be positive natural numbers with $2a\le b$. For a
restricted-weight input, define the deterministic restricted-weight query count to be
the product of its outer count and inner count.
\end{definition}

\begin{definition}[Deterministic query index pair]
\lean{FABL.deterministicRestrictedWeightQueryPair}
\label{FABL.deterministicRestrictedWeightQueryPair}
\leanok
\uses{FABL.RestrictedWeightInput, FABL.RestrictedWeightInput.innerInput, FABL.RestrictedWeightInput.outerInput, FABL.SmallBiasFourierInput.sampleCount, FABL.deterministicRestrictedWeightQueryCount}
Let $Q$ be the product of the outer and inner sample counts for a restricted weight
input. Each index $q	ext{ with }0\le q<Q$ determines the corresponding pair of indices,
one in the outer sample range and one in the inner sample range, under the standard
equivalence between a finite range of size equal to a product and the product of the two
ranges.
\end{definition}

\begin{definition}[Deterministic restricted weight query point]
\lean{FABL.deterministicRestrictedWeightQueryPoint}
\label{FABL.deterministicRestrictedWeightQueryPoint}
\leanok
\uses{FABL.RestrictedWeightInput, FABL.combineSignCube, FABL.deterministicRestrictedWeightQueryCount, FABL.deterministicRestrictedWeightQueryPair, FABL.restrictedWeightFixedAssignment, FABL.restrictedWeightFreeAssignment}
Given a restricted-weight input, a set $J$ of coordinates in $in{n}$, and a query index
$q$ from the set of all queries, the deterministic query point is the sign vector in
$\\{-1,1\\}^n$ obtained by combining the free assignment for the second component of the
query pair with the fixed assignment for its first component, using $J$.
\end{definition}

\begin{definition}[Estimate from deterministic query answers]
\lean{FABL.deterministicRestrictedWeightEstimateFromAnswers}
\label{FABL.deterministicRestrictedWeightEstimateFromAnswers}
\leanok
\uses{FABL.RestrictedWeightInput, FABL.RestrictedWeightInput.innerInput, FABL.RestrictedWeightInput.outerInput, FABL.Sign, FABL.SmallBiasFourierInput.sampleCount, FABL.deterministicRestrictedWeightQueryCount, FABL.indexedMonomial, FABL.restrictedWeightFreeAssignment, FABL.signValue}
Given a restricted-weight input, a finite coordinate set $J$, a subset $S\subseteq J$,
and a sign-valued answer for each valid query, define the estimate to be the average
over outer indices $o$ of the square of the inner average. The inner average is the
mean, over all inner indices $i$, of the answer at the query indexed by $(o,i)$
multiplied by the monomial indexed by $S$ evaluated at the free assignment associated
with $i$.
\end{definition}

\begin{theorem}[Restricted Fourier weight from oracle answers]
\lean{FABL.deterministicRestrictedWeightEstimateFromOracle_eq}
\label{FABL.deterministicRestrictedWeightEstimateFromOracle_eq}
\leanok
\uses{FABL.BooleanFunction, FABL.BooleanFunction.toReal, FABL.RestrictedWeightInput, FABL.deterministicRestrictedWeightEstimateFromAnswers, FABL.deterministicRestrictedWeightQueryPair, FABL.deterministicRestrictedWeightQueryPoint, FABL.restrictedFourierCoefficientEstimate, FABL.restrictedFourierWeightEstimate}
\difficulty{1}
Let the dimension and accuracy data satisfy $n>0$, $0<	ext{numerator}$, and $2	ext{
numerator}le	ext{denominator}$, and let the Fourier bound be at least $1$. For a
Boolean function $f$ on the sign cube, a coordinate set $J$, and a subset $S\subseteq
J$, the expected squared empirical Fourier coefficient obtained by querying $f$ at every
point of the prescribed product sample equals the restricted Fourier weight of $f$ on
$S$ over the fixed-coordinate sample.
\end{theorem}

\begin{definition}[Restricted weight local work]
\lean{FABL.deterministicRestrictedWeightLocalWork}
\label{FABL.deterministicRestrictedWeightLocalWork}
\leanok
\uses{FABL.RestrictedWeightInput, FABL.RestrictedWeightInput.innerCount, FABL.RestrictedWeightInput.outerCount}
For a restricted-weight input and a finite set $S$, the local work is
$N_{\mathrm{out}}\bigl(N_{\mathrm{in}}(|S|+2)+3\bigr)+1$, where $N_{\mathrm{out}}$ and
$N_{\mathrm{in}}$ are the outer and inner query counts associated with the input.
\end{definition}

\begin{definition}[Work of the restricted-weight construction]
\lean{FABL.deterministicRestrictedWeightConstructionWork}
\label{FABL.deterministicRestrictedWeightConstructionWork}
\leanok
\uses{FABL.RestrictedWeightInput, FABL.RestrictedWeightInput.innerInput, FABL.RestrictedWeightInput.outerInput, FABL.SmallBiasFourierInput.generatorInput, FABL.deterministicSmallBiasWork}
The deterministic restricted-weight construction work assigned to an input is the sum of
the deterministic small-bias construction work for its inner and outer generator inputs.
\end{definition}

\begin{definition}[Deterministic restricted-weight cost]
\lean{FABL.deterministicRestrictedWeightCost}
\label{FABL.deterministicRestrictedWeightCost}
\leanok
\uses{FABL.DeterministicQueryProgram.queryBatchCost, FABL.LearningCost, FABL.RestrictedWeightInput, FABL.deterministicRestrictedWeightConstructionWork, FABL.deterministicRestrictedWeightLocalWork, FABL.deterministicRestrictedWeightQueryCount}
Given a restricted-weight input and a finite set $S$, define the deterministic learning
cost as the componentwise sum of a cost with zero random-example calls and zero
membership queries whose work is the construction work, the deterministic query-batch
cost for the input, and a cost with zero random-example calls and zero membership
queries whose work is the local work for $S$.
\end{definition}

\begin{definition}[Deterministic restricted-weight program]
\lean{FABL.deterministicRestrictedWeightProgram}
\statementsource{Analysis-of-Boolean-Functions-by-Ryan-ODonnell}{
  pdf-ffbe9c7b3306-p157-b008,
  pdf-ffbe9c7b3306-p157-b009,
  pdf-ffbe9c7b3306-p158-b003
}
\label{FABL.deterministicRestrictedWeightProgram}
\leanok
\uses{FABL.DeterministicQueryProgram, FABL.RestrictedWeightInput, FABL.Sign, FABL.deterministicRestrictedWeightConstructionWork, FABL.deterministicRestrictedWeightEstimateFromAnswers, FABL.deterministicRestrictedWeightLocalWork, FABL.deterministicRestrictedWeightQueryCount, FABL.deterministicRestrictedWeightQueryPoint}
Given a restricted-weight input, a finite set $J\subseteq\mathrm{Fin}(n)$, and a subset
$S\subseteq J$, the deterministic query program first charges the construction work,
queries a batch of sign-cube points whose size is the specified query count, then
charges local work depending on the input and $S$, and returns the estimate determined
by the answers to that batch.
\end{definition}

\begin{theorem}[Output and cost of the restricted weight program]
\lean{FABL.DeterministicQueryProgram.runWithCost_deterministicRestrictedWeightProgram}
\label{FABL.DeterministicQueryProgram.runWithCost_deterministicRestrictedWeightProgram}
\leanok
\uses{FABL.BooleanFunction, FABL.DeterministicQueryProgram.runWithCost, FABL.RestrictedWeightInput, FABL.deterministicRestrictedWeightCost, FABL.deterministicRestrictedWeightEstimateFromAnswers, FABL.deterministicRestrictedWeightEstimateFromOracle_eq, FABL.deterministicRestrictedWeightProgram, FABL.deterministicRestrictedWeightQueryPoint, FABL.restrictedFourierWeightEstimate}
\difficulty{3}
Let $n$ be a positive integer, and let the input specify positive integers $a,d$ with
$2a\le d$ and an integer Fourier bound $s\ge1$. Let $J$ be any set of coordinates in the
$n$-dimensional sign cube, let $S\subseteq J$, and let $f$ be a Boolean function on that
cube. Running the deterministic restricted weight program with oracle $f$ returns the
restricted Fourier weight estimate for $f$ on $J,S$, together with a cost whose
random-example count is zero, whose membership-query count is the program's prescribed
query count, and whose work is the sum of its construction and local-work charges.
\end{theorem}

\begin{theorem}[Membership query count for restricted weight]
\lean{FABL.deterministicRestrictedWeightCost_queries}
\label{FABL.deterministicRestrictedWeightCost_queries}
\leanok
\uses{FABL.RestrictedWeightInput, FABL.RestrictedWeightInput.innerCount, FABL.RestrictedWeightInput.outerCount, FABL.deterministicRestrictedWeightCost}
\difficulty{3}
Let the input be a restricted-weight input with positive dimension, positive accuracy
numerator, accuracy at most $1/2$, and Fourier bound at least one. For any finite set
$S$, the membership-query component of the deterministic restricted-weight cost is the
product of the outer and inner sample counts.
\end{theorem}

\begin{theorem}[Random-example count for restricted weight]
\lean{FABL.deterministicRestrictedWeightCost_randomExamples}
\label{FABL.deterministicRestrictedWeightCost_randomExamples}
\leanok
\uses{FABL.RestrictedWeightInput, FABL.deterministicRestrictedWeightCost}
\difficulty{1}
Let the input be a restricted-weight input with positive dimension, positive accuracy
numerator, accuracy at most $1/2$, and Fourier bound at least one. For any finite set
$S$, the deterministic restricted-weight cost uses zero random examples.
\end{theorem}

\begin{theorem}[Work component of restricted-weight cost]
\lean{FABL.deterministicRestrictedWeightCost_work}
\label{FABL.deterministicRestrictedWeightCost_work}
\leanok
\uses{FABL.RestrictedWeightInput, FABL.deterministicRestrictedWeightConstructionWork, FABL.deterministicRestrictedWeightCost, FABL.deterministicRestrictedWeightLocalWork, FABL.deterministicRestrictedWeightQueryCount}
\difficulty{3}
For a restricted-weight input with positive dimension $n$, positive accuracy numerator
$a$, $2a$ at most its accuracy denominator, and Fourier bound $s \ge 1$, and for any
finite set $S$, the work component of its deterministic restricted-weight cost is the
sum of its construction work, query count, and local work for $S$.
\end{theorem}

\begin{definition}[Restricted weight algorithm scale]
\lean{FABL.RestrictedWeightInput.algorithmScale}
\label{FABL.RestrictedWeightInput.algorithmScale}
\leanok
\uses{FABL.RestrictedWeightInput, FABL.RestrictedWeightInput.innerInput, FABL.RestrictedWeightInput.outerInput, FABL.SmallBiasFourierInput.generatorInput, FABL.SmallBiasInput.scale}
For a restricted-weight input, the algorithm scale is the sum of the scales of its inner
and outer generator inputs.
\end{definition}

\begin{definition}[Restricted weight polynomial budget]
\lean{FABL.RestrictedWeightInput.polynomialBudget}
\label{FABL.RestrictedWeightInput.polynomialBudget}
\leanok
\uses{FABL.RestrictedWeightInput, FABL.RestrictedWeightInput.algorithmScale}
For a restricted-weight input with algorithm scale $A$, the polynomial budget is
$2^{20}(A+1)^8$.
\end{definition}

\begin{theorem}[Deterministic restricted weight cost resource bounds]
\lean{FABL.deterministicRestrictedWeightCost_resource_bounds}
\statementsource{Analysis-of-Boolean-Functions-by-Ryan-ODonnell}{pdf-ffbe9c7b3306-p158-b003}
\label{FABL.deterministicRestrictedWeightCost_resource_bounds}
\leanok
\uses{FABL.RestrictedWeightInput, FABL.RestrictedWeightInput.algorithmScale, FABL.RestrictedWeightInput.innerCount, FABL.RestrictedWeightInput.innerInput, FABL.RestrictedWeightInput.outerCount, FABL.RestrictedWeightInput.outerInput, FABL.RestrictedWeightInput.polynomialBudget, FABL.RestrictedWeightInput.quarterBias, FABL.SmallBiasFourierInput.generatorInput, FABL.SmallBiasFourierInput.sampleCount_le_four_mul_scale_sq, FABL.SmallBiasInput.n_le_scale, FABL.SmallBiasInput.polynomialBudget, FABL.SmallBiasInput.scale, FABL.deterministicRestrictedWeightConstructionWork, FABL.deterministicRestrictedWeightCost, FABL.deterministicRestrictedWeightCost_queries, FABL.deterministicRestrictedWeightCost_randomExamples, FABL.deterministicRestrictedWeightCost_work, FABL.deterministicRestrictedWeightLocalWork, FABL.deterministicRestrictedWeightQueryCount, FABL.deterministicSmallBiasWork, FABL.deterministicSmallBiasWork_le_polynomialBudget}
\difficulty{5}
% dataset-pipeline: stub -- prose pending build_informal_statements.py
Statement pending regeneration.
\end{theorem}

\begin{definition}[Restricted weight tasks]
\lean{FABL.RestrictedWeightTask}
\statementsource{Analysis-of-Boolean-Functions-by-Ryan-ODonnell}{pdf-ffbe9c7b3306-p158-b003}
\label{FABL.RestrictedWeightTask}
\leanok
\uses{FABL.RestrictedWeightInput}
A restricted-weight task consists of a restricted-weight input, a subset $J$ of the
finite index set determined by that input's dimension, and a subset $S\subseteq J$.
\end{definition}

\begin{definition}[Restricted weight task scale]
\lean{FABL.restrictedWeightTaskScale}
\label{FABL.restrictedWeightTaskScale}
\leanok
\uses{FABL.RestrictedWeightInput.algorithmScale, FABL.RestrictedWeightTask}
A restricted-weight task consists of a restricted-weight input, a finite coordinate set
$J$ in its dimension, and a subset of $J$. Its scale is the algorithm scale of its
input.
\end{definition}

\begin{theorem}[Deterministic restricted weight queries is big o]
\lean{FABL.deterministicRestrictedWeight_queries_isBigO}
\statementsource{Analysis-of-Boolean-Functions-by-Ryan-ODonnell}{pdf-ffbe9c7b3306-p158-b003}
\label{FABL.deterministicRestrictedWeight_queries_isBigO}
\leanok
\uses{FABL.RestrictedWeightTask, FABL.deterministicRestrictedWeightCost, FABL.deterministicRestrictedWeightCost_resource_bounds, FABL.restrictedWeightTaskScale}
\difficulty{3}
% dataset-pipeline: stub -- prose pending build_informal_statements.py
Statement pending regeneration.
\end{theorem}

\begin{theorem}[Deterministic restricted weight work is big o]
\lean{FABL.deterministicRestrictedWeight_work_isBigO}
\statementsource{Analysis-of-Boolean-Functions-by-Ryan-ODonnell}{pdf-ffbe9c7b3306-p158-b003}
\label{FABL.deterministicRestrictedWeight_work_isBigO}
\leanok
\uses{FABL.RestrictedWeightTask, FABL.deterministicRestrictedWeightCost, FABL.deterministicRestrictedWeightCost_resource_bounds, FABL.restrictedWeightTaskScale}
\difficulty{3}
% dataset-pipeline: stub -- prose pending build_informal_statements.py
Statement pending regeneration.
\end{theorem}
